\ifja
    \chapter{魔術の残骸と、冷徹なる計算機}
    \NLP（神経言語プログラミング）が「再現性のない疑似科学」と断じられて久しい。アカデミアの書棚から追放され、自己啓発セミナーの商材として消費される現状を見て、「やはりインチキだったのか」と失望した者も多いだろう。あるいは、ラポールやイエスセットといった表層的なテクニックの一片を切り取り、「効果がある」と安易に信じ込んでいる者もいる。

    だが、現場の真実を知る者は知っている。
    私たちは、\jpDrBandler が放つ圧倒的な「魔術」をこの目で目撃してきた。重度の恐怖症が一瞬で霧散し、数十年来の頑なな信念がわずか数分の対話で解体され、再構築される。彼の指先ひとつ、声のトーンひとつで、人間の知覚空間そのものが歪められる。彼は間違いなく本物である。

    しかし――なぜ私たちは、彼と同じスクリプトを読み上げ、同じ催眠言語パターンをなぞっても、一向に相手をトランスに落とせないのか。なぜ他人はこちらの指示に従わず、強固な現実の壁に跳ね返されるのか。

    答えは残酷なほどシンプルだ。
    私たちは彼が「見せている手品（表層の言語）」を表面的に真似ていただけで、彼がハックしていた「生体情報処理ハードウェアの物理法則」を何ひとつ理解していなかったからである。

    \NLPを神秘のベールから引き剥がし、真の工学として掌中に収めるためには、私たちが「心」や「無意識」に対して抱いてきた自明の思い込みはもちろんのこと、\NLPという体系そのものが掲げる前提すら疑う冷徹な批判的知性が必要となる。

    種明かしを始める前に、まずはこの生体コンピュータの動作原理を記述するための、最低限の数学的・情報論的プロトコルを定義しておこう。
\else

    \chapter{The Debris of Magic and the Cold Computer}

    It has long been asserted that \NLP is an unreproducible pseudoscience. Cast out from academic bookshelves and consumed as mere commodity in self-help seminars, it has left many disappointed, concluding that it was nothing more than a fraud all along. Conversely, others cling naively to superficial fragments—rapport building, yes-sets—convinced of their efficacy without scrutiny.

    Yet those who know the reality of the field know better.
    We have witnessed firsthand the overwhelming "magic" wielded by \exDrBandler. Severe phobias dissolve into thin air in an instant; entrenched beliefs held for decades are dismantled and reconstructed in a matter of minutes. With a mere flick of his finger or a subtle shift in vocal tone, the subject's perceptual manifold itself bends. He is, beyond any shadow of a doubt, the real thing.

    And yet—why is it that when we recite his exact scripts and trace his hypnotic language patterns verbatim, our subjects remain stubbornly untranced? Why do others resist our directions, rebounding off the unyielding wall of consensus reality?

    The answer is brutally simple.
    We were merely mimicking the surface sleight-of-hand—the overt syntax—while remaining completely blind to the physical laws of the biological information-processing hardware he was actually hacking.

    To strip \NLP of its mystical veil and claim it as a true discipline of engineering demands a cold, critical intellect—one willing to dismantle not only our reflexive assumptions regarding the "mind" and the "unconscious," but also the foundational dogmas postulated by the framework of \NLP itself.

    Before we reveal the machinery behind the illusion, let us first formalize the baseline mathematical and information-theoretic protocols required to describe the operational principles of this biological computer.
\fi

\ifja
    \section{ベイズ統計学の基礎数理}
    ベイズ統計学とは、新しいデータや証拠（エビデンス）を獲得するたびに、自らの見立て（確信度）を柔軟に更新していく実践的な確率論である。
    無限回の繰り返し試行を前提とする従来の頻度主義統計学とは異なり、ベイズ統計学は「一度きりの事象」に対しても主観確率を割り振って評価できる。これにより、情報が極めて乏しい初期段階であっても事前知識をもとに暫定的な推論を立ち上げることができる。
\else
    \section{Foundational Mathematics of Bayesian Statistics}
    Bayesian statistics is a pragmatic branch of probability theory wherein one iteratively updates their conviction (degree of belief) upon acquiring new data or empirical evidence.
    Unlike traditional frequentist statistics, which relies on the presupposition of infinitely repeatable trials, Bayesian statistics assigns and evaluates subjective probabilities even for single, non-repeatable events. This capability enables the initialization of tentative inferences derived from prior knowledge, even in baseline regimes where observational data is exceptionally sparse.
\fi

\ifja
    \subsection{事象の確率と\jpTextConditionalProbability }
    ベイズ推論を理解するための基礎となるのが「\jpTextConditionalProbability（\exTextConditionalProbability）」である。

    確率が正である事象 $\eventB$（すなわち $\Probability{\eventB} > 0$）がすでに観測された条件下で、事象 $\eventA$ が生起する条件付き確率は次式で定義される。

    \begin{equation}
    \ConditionalProbability{\eventA}{\eventB} = \frac{\Probability{\eventA \cap \eventB}}{\Probability{\eventB}}
\label{eq:cond_prob}
\end{equation}




    ここで $\Probability{\eventA \cap \eventB}$ は事象 $\eventA$ と事象 $\eventB$ が同時に生起する同時確率を表す。この定義式を変形すると、確率の乗法定理が得られる。

    \begin{equation}
    \Probability{\eventA \cap \eventB} = \ConditionalProbability{A}{B} \Probability{\eventB} = \ConditionalProbability{\eventB}{\eventA} \Probability{\eventA}
\label{eq:mult_rule}
\end{equation}



    \eqref{eq:mult_rule} 式は、事象 $\eventA$ と $\eventB$ が同時に生起するプロセスを「$\eventB$ を前提とした $\eventA$ の生起」と捉えても、「$\eventA$ を前提とした $\eventB$ の生起」と捉えても数学的に同値であることを示している。
\else
    \subsection{Probability of Events and \exTextConditionalProbability}
    The foundation for understanding Bayesian inference is ``\exTextConditionalProbability.''

    Given the condition that an event $\eventB$ with positive probability (i.e., $\Probability{\eventB} > 0$) has already been observed, the conditional probability that event $\eventA$ occurs is defined by the following equation:

    \begin{equation}
    \ConditionalProbability{\eventA}{\eventB} = \frac{\Probability{\eventA \cap \eventB}}{\Probability{\eventB}}
\label{eq:cond_prob}
\end{equation}




    where $\Probability{\eventA \cap \eventB}$ denotes the joint probability that events $\eventA$ and $\eventB$ occur simultaneously. Rearranging this definition yields the multiplication rule of probability:

    \begin{equation}
    \Probability{\eventA \cap \eventB} = \ConditionalProbability{A}{B} \Probability{\eventB} = \ConditionalProbability{\eventB}{\eventA} \Probability{\eventA}
\label{eq:mult_rule}
\end{equation}



    Equation~\eqref{eq:mult_rule} indicates that the process in which events $\eventA$ and $\eventB$ co-occur is mathematically equivalent whether interpreted as ``the occurrence of $\eventA$ given $\eventB$'' or ``the occurrence of $\eventB$ given $\eventA$.''
\fi

\ifja
    \subsection{\jpTextBayesTheoremの導出}
    乗法定理 \eqref{eq:mult_rule} の両辺を $\Probability{\eventB}$（ただし $\Probability{\eventB} > 0$）で除すことにより、基礎的な\jpTextBayesTheorem（\exTextBayesTheorem）が導かれる。
    \begin{equation}
    \ConditionalProbability{\eventA}{\eventB} = \frac{\ConditionalProbability{\eventB}{\eventA} \Probability{\eventA}}{\Probability{\eventB}}
\label{eq:bayes_basic}
\end{equation}



\else
    \subsection{Derivation of \exTextBayesTheorem}
    Dividing both sides of the multiplication rule \eqref{eq:mult_rule} by $\Probability{\eventB}$ (where $\Probability{\eventB} > 0$) yields the elementary form of \exTextBayesTheorem:
    \begin{equation}
    \ConditionalProbability{\eventA}{\eventB} = \frac{\ConditionalProbability{\eventB}{\eventA} \Probability{\eventA}}{\Probability{\eventB}}
\label{eq:bayes_basic}
\end{equation}



\fi

\subsection{全確率の定理を用いた一般化}
\subsection{ベイズ更新（逐次更新）のメカニズム}
\section{モンティ・ホール問題：人間の直感の破綻}
\subsection{全米を巻き込んだ大論争}
\subsection{ベイズの定理による厳密解}
\section{スコーピオン号の沈没とベイズ探索理論：不確実性下における確率空間の収束}
\subsection{深海の確率格子とモンテカルロ・サンプリング}
\subsection{直感を排除した「客観的確信度」の到達点}
\section{人間の神経系と数学的ベイズの違い}


